% ============================================================
%  Chapter scaffold -- theory-reference skill
%  Replace CHAPTER_TITLE, LABEL, and all section content.
%  Do not use \newpage or \clearpage anywhere in this file.
% ============================================================

\chapter{CHAPTER\_TITLE}
\label{ch:LABEL}

% ── Roadmap (required, once per chapter) ────────────────────
\begin{roadmap}
\textbf{What we build.}
ONE_SENTENCE_SUMMARY_OF_WHAT_THIS_CHAPTER_ESTABLISHES.

After this chapter you will be able to:
\begin{itemize}[noitemsep]
  \item OBJECTIVE_1
  \item OBJECTIVE_2
  \item OBJECTIVE_3
\end{itemize}

\textbf{Prerequisites:} PREREQUISITES.
\end{roadmap}

% ============================================================
\section{SECTION\_TITLE}
\label{sec:LABEL_sec1}
% ============================================================

ONE_MOTIVATING_SENTENCE_BEFORE_ANY_DEFINITION.

\begin{definition}[DEFINITION_NAME]
\label{def:LABEL}
DEFINITION_BODY.
\end{definition}

% Use keyidea when the conceptual insight is non-obvious
\begin{keyidea}
KEY_INSIGHT.
\end{keyidea}

\begin{theorem}[THEOREM_NAME]
\label{thm:LABEL}
THEOREM_STATEMENT.
\end{theorem}

\begin{proof}
PROOF_BODY.
\end{proof}

% Use keyeqn for the 1-2 most important equations in the section
\begin{keyeqn}
\begin{equation}
  IMPORTANT\_EQUATION
  \label{eq:LABEL}
\end{equation}
SHORT_ANNOTATION.
\end{keyeqn}

\begin{checkpoint}
VERIFICATION_QUESTION.
\end{checkpoint}

\begin{refnote}
CITATION, Chapter/Section REFERENCE.
\end{refnote}

% ============================================================
\section{SECTION\_TITLE\_2}
\label{sec:LABEL_sec2}
% ============================================================

% ... continue pattern above

% ── Optional: connection box near end of chapter ────────────
% \begin{connection}{DOMAIN Connection}
% HOW_THIS_CHAPTER_CONNECTS_TO_THE_TARGET_APPLICATION.
% \end{connection}
